\chapter{Deficiency, excess, and interpretation}
\section{Learning objectives}
Construct a differential diagnosis for suspected nutrient deficiency; choose investigations rationally; recognize toxicity and poisoning; and design a monitored replacement plan that addresses the cause.

\section{Why deficiency happens}
Inadequacy may result from low intake, food insecurity, restrictive eating, alcoholism, increased physiological need, impaired digestion or absorption, kidney disease, chronic inflammation, blood loss, or medication effects. A deficiency is not proven by a nonspecific symptom. Fatigue, brittle nails, cramps, poor concentration, and hair changes have many causes.

\section{A clinical pattern}
Evaluation usually begins with history: diet, weight change, gastrointestinal symptoms, menstrual or other blood loss, sun exposure, medicines, alcohol, and supplements. A clinician may order a blood count, metabolic tests, or nutrient-specific tests. Treatment should address the cause as well as replenish the nutrient. Follow-up matters because some deficiencies take time to correct and some supplements alter test results.

\section{Recognizing excess}
Toxicity can result from supplements, fortified foods, therapeutic medicines, occupational exposure, or inherited disorders of mineral handling. The history should establish the exact product, elemental dose, number of units, timing, and all concurrent products. Symptoms may be delayed, and the absence of symptoms does not establish safety.
Excess may be acute, as with accidental ingestion, or chronic, as with repeated megadoses. Possible signals include nausea, severe constipation or diarrhea, weakness, confusion, abnormal bleeding, jaundice, kidney-stone symptoms, or tingling. Fat-soluble vitamins and minerals such as iron, selenium, and iodine deserve particular respect, but water-soluble vitamins can also be toxic in excess.

\section{When urgent help is appropriate}
Seek urgent poison-control or emergency advice after a child swallows iron or another concentrated supplement, after a very large intentional dose, or when there is persistent vomiting, confusion, collapse, seizures, breathing difficulty, severe weakness, or chest pain. Keep the container and estimate the amount and time taken.

\warning{Do not induce vomiting unless poison-control personnel specifically instruct you to do so. Emergency numbers and poison-control services differ by country.}
